\documentclass[11pt]{article}
\usepackage[margin=1in]{geometry}
\usepackage{amsmath,amssymb,amsthm,enumitem}
\DeclareMathOperator{\nul}{null}
\DeclareMathOperator{\range}{range}
\newcommand{\F}{\mathbf{F}}
\newcommand{\C}{\mathbf{C}}
\newcommand{\LL}{\mathcal{L}}
\newcommand{\prob}[2]{\medskip\noindent\textbf{#1.}\ \textnormal{(#2)}\ }
\setlength{\parindent}{0pt}
\setlength{\parskip}{6pt}
\pagestyle{empty}
\begin{document}
\begin{center}
{\Large MATH 3406 Exam 2: Extra Practice}\\[4pt]
{\small Twelve exercises from Axler, \emph{Linear Algebra Done Right} (4th ed.), in decreasing order of difficulty}
\end{center}

Throughout, $V$ and $W$ are vector spaces over $\F$. Nothing is assumed finite-dimensional unless stated.

\prob{1}{5B \#4} Suppose $\F=\C$, $T\in\LL(V)$, $p\in\mathcal{P}(\C)$ is a nonconstant polynomial, and $\alpha\in\C$. Prove that $\alpha$ is an eigenvalue of $p(T)$ if and only if $\alpha=p(\lambda)$ for some eigenvalue $\lambda$ of $T$.

\textbf{Forward}: Suppose that $\alpha$ is an eigenvalue of $p(T).$ Then, we know 
$$
\dim \ker (p(T) - \alpha I) > 0,
$$
i.e., $p(T) - \alpha I$ is not injective. Factor $p(z) - \alpha$ as
$$
c(z - \mu_1)(z - \mu_2) \ldots (z - \mu_m),
$$
yielding
$$
c(T- \mu_1 I ) \ldots (T - \mu_m I ) v = 0.
$$
for some $v \neq 0.$ Since this map is not injective, there must be some 
$$
T - \mu_k I
$$
that is not injective, meaning such a $\mu_k = \lambda$ for some eigenvalue $\lambda.$ Then, $p(\lambda) - \alpha$


\textbf{Backward}: Conversely, know that $T$ is invariant under $\ker T - \lambda I,$ for any eigenvalue $\lambda$ of $T.$ Hence, for any $v \in \ker T - \lambda I,$ we get 
$$
p(T) v = p(\lambda) v,
$$
meaning $p(\lambda)$ is an eigenvalue of $p(T).$ Since $\alpha = p(\lambda),$ we know $\alpha$ is an eigenvalue of $p(T).$

\prob{2}{3F \#33} Suppose $U$ is a subspace of $V$, and let $\pi\colon V\to V/U$ be the quotient map, so $\pi'\in\LL\bigl((V/U)',V'\bigr)$.
\begin{enumerate}[label=(\alph*),topsep=2pt,itemsep=0pt]
\item Show that $\pi'$ is injective.
\item Show that $\range \pi'=U^0$.
\item Conclude that $\pi'$ is an isomorphism from $(V/U)'$ onto $U^0$.
\end{enumerate}

\prob{3}{3D \#6} Suppose $W$ is finite-dimensional and $S,T\in\LL(V,W)$. Prove that $\nul S=\nul T$ if and only if there exists an invertible $E\in\LL(W)$ such that $S=ET$.

\prob{4}{3F \#22} Suppose $V$ is finite-dimensional and $U$, $W$ are subspaces of $V$.
\begin{enumerate}[label=(\alph*),topsep=2pt,itemsep=0pt]
\item Show that $(U+W)^0=U^0\cap W^0$.
\item Show that $(U\cap W)^0=U^0+W^0$.
\end{enumerate}

\prob{5}{5C \#12a} Suppose $V$ is finite-dimensional, $T\in\LL(V)$ has an upper-triangular matrix with respect to some basis of $V$, and $U$ is a subspace of $V$ invariant under $T$. Prove that $T|_U$ has an upper-triangular matrix with respect to some basis of $U$.

\prob{6}{5B \#28} Suppose $V$ is finite-dimensional and $T\in\LL(V)$. Prove that the minimal polynomial of $T'\in\LL(V')$ equals the minimal polynomial of $T$.

\prob{7}{5D \#8} Suppose $V$ is finite-dimensional and $T\in\LL(V)$. Let $\lambda_1,\dots,\lambda_m$ be the distinct nonzero eigenvalues of $T$. Prove that
\[\dim E(\lambda_1,T)+\cdots+\dim E(\lambda_m,T)\le\dim\range T.\]

\prob{8}{3F \#21} Suppose $V$ is finite-dimensional and $U$, $W$ are subspaces of $V$.
\begin{enumerate}[label=(\alph*),topsep=2pt,itemsep=0pt]
\item Prove that $W^0\subseteq U^0$ if and only if $U\subseteq W$.
\item Prove that $W^0=U^0$ if and only if $U=W$.
\end{enumerate}

\prob{9}{5A \#25} Suppose $T\in\LL(V)$ and $u$, $w$ are eigenvectors of $T$ such that $u+w$ is also an eigenvector of $T$. Prove that $u$ and $w$ correspond to the same eigenvalue.

\prob{10}{3E \#6} Suppose $v,x\in V$ and $U$, $W$ are subspaces of $V$ such that $v+U=x+W$. Prove that $U=W$.

\prob{11}{3F \#3} Suppose $V$ is finite-dimensional and $v\in V$ with $v\ne0$. Prove that there exists $\varphi\in V'$ such that $\varphi(v)=1$.

\prob{12}{5A \#30} Suppose $T\in\LL(V)$ and
\[(T-2I)(T-3I)(T-4I)=0.\]
Suppose $\lambda$ is an eigenvalue of $T$. Prove that $\lambda=2$ or $\lambda=3$ or $\lambda=4$.
\end{document}